%% kapa: exam=1 round=29 year=2018 period=1 form=A subject=economics subject_ko=경제학원론 seq=20 number=60
%% source: Q-Net 감정평가사 기출문제 (cst00309) artlSeq=5202808 file=extracted/1008/5202808/2018년 제29회 감정평가사 1차시험 기출문제/2018년 제29회 감정평가사 1차시험 기출문제/1교시 A형.pdf
%% conversion: kordoc 4.24.51 -> tools/gen1.py (auto); review=proofread-sonnet; flags=none
\begin{question}{60}
( )에 들어갈 내용으로 옳은 것은?
\begin{framedtext}
○ 생산의 긍정적 외부효과가 있을 때, ( ㄱ )이 ( ㄴ )보다 작다. \par
○ 소비의 부정적 외부효과가 있을 때, ( ㄷ )이 ( ㄹ )보다 크다. \par
\end{framedtext}
\begin{choices}
  \item ㄱ: 사회적 한계비용, ㄴ: 사적 한계비용, ㄷ: 사회적 한계편익, ㄹ: 사적 한계편익
  \item ㄱ: 사회적 한계비용, ㄴ: 사적 한계비용, ㄷ: 사적 한계편익, ㄹ: 사회적 한계편익
  \item ㄱ: 사적 한계비용, ㄴ: 사회적 한계비용, ㄷ: 사회적 한계편익, ㄹ: 사적 한계편익
  \item ㄱ: 사적 한계비용, ㄴ: 사회적 한계비용, ㄷ: 사적 한계편익, ㄹ: 사회적 한계편익
  \item ㄱ: 사회적 한계편익, ㄴ: 사적 한계편익, ㄷ: 사적 한계비용, ㄹ: 사회적 한계비용
\end{choices}
\end{question}
\answer{2}
